% Comparación de estrategias, riesgo residual, liquidez (colateral), tasa (no recomendada hoy), condiciones y límites.
\section{Análisis comparativo e implementación}

Ambas estrategias reducen el riesgo en la magnitud que importa y no exigen pago inicial. Lo que puede impedir
ejecutarlas no es su costo, sino el colateral que pediría el banco. El Cuadro~\ref{tab:comparacion} las compara con
los mismos criterios.

\begin{table}[htbp]\centering\small
\caption{Comparación: E1 y E2 reducen el riesgo sin costo inicial; el límite es el colateral}\label{tab:comparacion}
\renewcommand{\arraystretch}{1.25}
\newcommand{\filasep}{\arrayrulecolor{gray!45}\specialrule{0.3pt}{1pt}{1pt}\arrayrulecolor{black}}
\newcommand{\encab}[2]{\textbf{#1}\newline #2}
\begin{tabularx}{\linewidth}{>{\raggedright\arraybackslash\bfseries}p{2.3cm}
    >{\raggedright\arraybackslash}X
    >{\raggedright\arraybackslash}X
    >{\raggedright\arraybackslash}X}
\toprule
\normalfont\textbf{Criterio}
  & \cellcolor{prioAlta}\encab{E1: tipo de cambio}{CCS + NDF}
  & \cellcolor{prioAlta}\encab{E2: precio del crudo}{Swap + collar WTI}
  & \cellcolor{prioMedia}\encab{Tasa de interés}{\textit{No recomendada hoy}} \\
\midrule
Exposición cubierta & \CoberturaFX\ de la posición en S/
  & \InvCubiertoMMbl\ MMbl, \CoberturaCrudoPct\ del inventario \newline (US\$~\InvCubiertoValor\ MM)
  & Tasa base de la refinanciación del CESCE (US\$~\CESCESaldo\ MM) \\ \filasep
Reducción del riesgo \newline {\normalfont (US\$ MM)}
  & Pérdida si el TC cae 10\,\%: \newline Sin cobertura: (\EscSinCobMenosDiez) \newline Con cobertura: (\EscConCobMenosDiez)
  & Pérdida si el WTI cae 30\,\%: \newline Sin cobertura: (\PerdidaSinCobTreinta) \newline Con cobertura: (\PerdidaMaxMixto)
  & Baja: no cubre el spread de crédito \\ \filasep
Costo inicial & Cero; CVA de \CCSCVApb~pb en la tasa & Cero (swap y collar) & Cero (swap de inicio diferido) \\ \filasep
Costo económico & Diferencial implícito del NDF: \NDFCostoAnual\ anual
  & Swap: US\$~\CostoSwap\ MM frente al spot (nulo frente al forward)
  & Pagos si la tasa baja \\ \filasep
Colateral con umbral cero \newline {\normalfont (US\$ MM)} & \ColateralTC\ si el TC sube 10\,\% & \ColateralWTI\ si el WTI sube 30\,\% & Bajo \\ \filasep
Tratamiento NIIF~9 & CCS: flujos de efectivo (ORI); NDF: valor razonable en resultados
  & Valor razonable del inventario existente; se redesigna al rotar
  & Flujos de efectivo, solo si la refinanciación es altamente probable \\
\bottomrule
\end{tabularx}
\fuente{Cuadros~\ref{tab:e1-param} a \ref{tab:e2-escenarios}; \textcite{ifrs9}. Elaboración propia. El inventario cubierto se valora al WTI de la fecha de valorización; al cierre de 2025 el inventario de crudo valía US\$~\InvCrudoUSD\ MM (Cuadro~\ref{tab:matriz}).}
\end{table}

\textbf{Riesgo residual.} El VaR al 95\,\% es la pérdida que solo se supera en uno de cada veinte casos. Al
vencimiento de cada cobertura, baja de US\$~\VaRFXSin\ a \VaRFXCon\ millones en la E1 y de \VaRCrudoSin\ a
\VaRCrudoCon\ millones en la E2 (Cuadro~\ref{tab:var}). La E2 conserva más riesgo porque deja abierto el 20\,\% del
inventario y la zona entre los strikes del collar.

\textbf{Liquidez: la restricción que decide.} Los derivados OTC no tienen márgenes de bolsa, pero el banco pedirá
colateral a una contraparte B-/Caa1. El \emph{umbral} del anexo CSA es la pérdida que el banco tolera antes de pedir
garantías en efectivo; se supone umbral cero, el caso más exigente. Un choque instantáneo conjunto (TC $+10$\,\% y
WTI $+30$\,\%) exigiría US\$~\ColateralConjunto\ millones, pero el CCS vive 27 meses: lo relevante es la
exposición potencial futura (PFE) al 99\,\% en toda su vida. Con el NDF y la E2 (correlación TC--WTI de \CorrTCWTI)
llega a US\$~\PFENoventaNueve\ millones, frente a US\$~\LineasLibres\ millones de líneas libres
(Cuadro~\ref{tab:colateral}); el préstamo puente (US\$~\PrestamoPuente\ millones) tiene otro destino. La ejecución
exige un umbral total de al menos US\$~\UmbralNecesario\ millones (unos \UmbralPorBanco\ por banco, con
\NContrapartes), que deja libre el \ColchonLineas\ de las líneas, o la garantía del Estado; sin ello, la cobertura
agravaría el riesgo de liquidez que busca mitigar \parencite{moodys2026petroperu}.

\medskip\noindent
\begin{minipage}[t]{0.42\linewidth}\footnotesize\setlength{\tabcolsep}{2.5pt}% ARCHIVO GENERADO — NO EDITAR A MANO
\centering
\captionof{table}{VaR al 95\,\% al vencimiento de la cobertura (US\$ MM)}\label{tab:var}
\begin{tabular}{lrrr}
\toprule
 & \textbf{Sin cob.} & \textbf{Con cob.} & \textbf{Reducción} \\
\midrule
E1 (TC) & 78.5 & 3.1 & 96\% \\
E2 (crudo) & 93.3 & 36.5 & 61\% \\
\bottomrule
\end{tabular}
\fuente{BCRP, series estadísticas; NYMEX CL vía Yahoo Finance; CBOE OVX vía FRED (28-sep-2026). Volatilidad realizada de un año. Elaboración propia. Horizonte: 91 días (E1) y vencimiento de las opciones (E2). La E2 incluye el 20 \% no cubierto.}
\end{minipage}\hfill
\begin{minipage}[t]{0.55\linewidth}\footnotesize\setlength{\tabcolsep}{2.5pt}% ARCHIVO GENERADO — NO EDITAR A MANO
\centering
\captionof{table}{Colateral exigible con umbral cero (US\$ MM)}\label{tab:colateral}
\begin{tabular}{lrrrr}
\toprule
 & \textbf{E1} & \textbf{E2} & \textbf{Colat.} & \textbf{Déficit} \\
\midrule
TC $+10$\,\% & 121.6 & 0.0 & 121.6 & 99.4 \\
WTI $+30$\,\% & 0.4 & 51.8 & 52.2 & 30.0 \\
TC $+10$\,\% y WTI $+30$\,\% & 121.6 & 51.8 & 173.4 & 151.2 \\
PFE 95\,\% & 108.3 & 97.5 & 168.9 & 146.7 \\
PFE 99\,\% & 151.4 & 157.8 & 253.6 & 231.4 \\
\bottomrule
\end{tabular}
\fuente{\textcite{petroperu2026eeff}, Nota 3.1(c); BCRP, series estadísticas; NYMEX CL vía Yahoo Finance; CBOE OVX vía FRED (28-sep-2026). Elaboración propia. E1 y E2: pasivo de los derivados, a TC medio. Choques: instantáneos. PFE: pico en la vida del CCS más NDF a 91 días (E1) y programa rodante a su plazo (E2), lognormal con volatilidad realizada; E1 y E2 agregadas con correlación TC--WTI de 0.35. Déficit: colateral menos líneas libres (US\$ 22.2 MM).}
\end{minipage}
\medskip

\textbf{Plan B si el banco no concede el umbral.} (a) Escalonar: contratar primero el 50\,\% del CCS, con lo que el
colateral ante un TC $+10$\,\% baja a US\$~\ColateralTCMitad\ millones. (b) Comprar opciones, que solo exigen prima y
ningún colateral; el put de WTI sobre el \CoberturaCrudoPct\ cuesta US\$~\CostoPutTotal\ millones, por lo que se
limitaría al tramo más expuesto. (c) Pactar un umbral escalonado según la calificación y un reajuste semanal del
colateral, en lugar de diario. (d) Obtener la garantía del Estado sobre la línea de derivados.

\textbf{Riesgo de tasa: no se recomienda cubrirlo hoy.} La deuda de largo plazo ya es fija; pasarla a variable añadiría riesgo sin
bajar el costo de crédito, que lo domina el spread. Queda como opción un swap de inicio diferido para la
refinanciación del CESCE (tasa de \FwdSwapRate\ con pago anual), solo cuando esa refinanciación sea altamente probable
(NIIF~9, 6.3.3). Hoy no lo es, por el incumplimiento del covenant.

\textbf{Condiciones de implementación.} (i) El Directorio aprueba una Política de Cobertura con límites (\LimiteFX\
de la posición en soles y \LimiteCrudo\ del inventario de crudo), prohibición expresa de especular y métricas
mensuales. (ii) Se firman contratos ISDA con anexo CSA y un umbral de al menos US\$~\UmbralNecesario\ millones, o se
aplica el plan B, con período de reposición del colateral y monto mínimo de transferencia. Como el covenant del
CESCE ya está incumplido, el incumplimiento cruzado del ISDA (§5(a)(vi)) se limita a deudas efectivamente aceleradas
(\emph{cross-acceleration}), con umbral alto y dispensa expresa para el CESCE; sin ella, no se firma. Se negocia además que la rebaja de calificación, la opinión con salvedades o la incertidumbre de
empresa en marcha no causen el cierre inmediato sin un período de cura. También una cláusula de prepago: si el BN
prepaga el préstamo, el CCS debe poder cancelarse a valor de mercado. (iii) Antes de pactar, la Gerencia Legal opina y
se obtiene la autorización que exija el régimen de endeudamiento y tesorería de una empresa del Estado. Se evalúa si
los compromisos contingentes que el D.U.~\DUCrisis\ autoriza al MINEM (hasta US\$~\DUCrisisMonto\ millones, Nota~1,
p.~25) pueden respaldar las líneas de derivados. (iv) Las coberturas se documentan según la NIIF~9 y se revelan según la NIIF~7, ¶21A--24G
\parencite{ifrs9}; los swaps con Citibank entran a la política. (v) Se separan funciones: Finanzas contrata; Riesgos controla límites y verifica de forma independiente los
precios de los bancos; Contabilidad registra \parencite{humala2026}. La política fija límites por contraparte y de pérdida, y Riesgos
aprueba los modelos; con insumos proxy y el CVA, el nivel de valor razonable (2 o 3; NIIF~13, ¶73) se determina por
instrumento. El reporte
regulatorio se coordina con cada banco.

\textbf{Supuestos que limitan el resultado.} Las cifras son indicativas. El cronograma del BN es inferido (tasa
del CCS entre \CCSTasaMin\ y \CCSTasaMax). La posición en soles distinta del BN y el inventario son de diciembre de
2025: si cambiaron, cambian los nocionales (nunca por encima de la exposición al pactar), no el diseño. La curva soberana y el OVX sustituyen a la curva swap y a la
superficie de volatilidad, que no son públicas. Antes de pactar se requieren el cronograma contractual, la posición
actualizada y cotizaciones ejecutables de al menos tres bancos (no BlueGamma ni Yahoo Finance).
